\newpage
{\color{gray}\hrule}
\begin{center}
\section{Conclusioni}
\bigskip
\end{center}
{\color{gray}\hrule}
\vspace{0.5cm}
Nella prima parte dell'esperienza si è caratterizzato il comportamento di un \textit{Silicon Photomultiplier}, stimandone in particolare il \textit{dark count rate} e l'\textit{optical crosstalk}: si è anche verificato (seppur con dei valori di $\chi^2$ bassi) che il DCR segue un andamento lineare crescente al variare della temperatura di funzionamento, mentre l'OCT rimane pressochè costante.

Utilizzando una sorgente LED, si è prima stimato il guadagno del SiPM e si è poi studiata la risposta dell'apparato di misura, composto da una \textit{Power Supply and Amplification Unit} e da un \textit{Digitizer}, al variare di diverse impostazioni del setup sperimentale: con la relazione lineare tra il \textit{gain} digitale $G_{ampl}$ e la distanza tra i picchi $\Delta_{pp}$, si è dimostrato che l'apparato risponde linearmente alla quantità di carica restituita dall'accensione delle celle del SiPM; mentre con la relazione lineare tra la tensione di alimentazione $V_{bias}$ e il guadagno del SiPM si è verificato che la risposta del SiPM all'aumentare della tensione di polarizzazione inversa delle celle è lineare. Si è infine stimato il numero medio di celle accese al variare dell'intensità del LED in modo approssimativo, mentre non è stato possibile stimarlo con il metodo più accurato.

Nella seconda parte dell'esperienza si è utilizzato il SiPM per studiare il decadimento $\beta^-$ dello ${}^{90}$Sr: si è estrapolato con la tecnica del plot di Kurie l'\textit{end point} $T_{0,Y}$ dello spettro per stimare il coefficiente di conversione tra le misure in CHN e il loro valore in energia (sia per sorgenti allineate sia per sorgenti non allineate al SiPM) e lo si è usato per stimare il massimo della distribuzione $T_{max,\text{Sr}}$. Si è poi studiato l'assorbimento della radiazione $\beta^-$ in foglietti di carta e alluminio, ricavando da un lato il valore del coefficiente di assorbimento della carta, il quale risulta nel range di valore atteso, e dall'altro stimando lo spessore dei foglietti di alluminio (di difficile misura diretta) conoscendone il coefficiente di assorbimento.

%paolo dimmi se ti sembra grammaticalmente corretto e se ha senso 
Si è poi considerato lo spettro del ${}^{137}$Cs, dove tramite l'analisi dei picchi a bassi canali si è stimata la risoluzione attesa dell'apparato, la quale risulta statisticamente compatibile con quella ricavata dal picco di conversione interna. Si è infine stimato il valore massimo della distribuzione $T_{max,\text{Cs}}$, che non risulta compatibile con il valore atteso per una sottostima dell'errore.

Con l'utilizzo della \textit{Cosmic Box} è stato possibile acquisire uno spettro e stimare l'energia in media depositata dai muoni nello scintillatore. 
Tale valore è risultato statisticamente compatibile con la previsione teorica della legge di Bethe-Bloch. La risoluzione dell'apparato ottenuta dallo spettro dei muoni risulta maggiore di quella ottenuta dallo studio dello spettro del ${}^{137}$Cs. Si attribuisce questo aumento in primo luogo al piccolo numero di eventi registrati e in secondo luogo al fatto che è stato possibile effettuare una regressione solo sulla coda sinistra dei dati e non sulla coda destra poiché questa è influenzata dai muoni fuori asse.
